\section{问题三的模型与求解}
\label{sec:q3}

\subsection{从单源定位到搜索与清除协同}

第三问还需在源数未知时完成搜索与清除。逐源处理容易绕行，先搜完再清除又难以及时调整行程。本文采用\textbf{基于可行区域的搜索—定位—清除协同策略，定稿版本为 v6}：共同安排覆盖扫描与目标处理，并据反馈更新路线。

Vander Hook等将移动与测量时间共同纳入主动定位评价\cite{VanderHook2014}。结合本题计费规则，记总路程为$D_{\mathrm{move}}$，无线电检测、实际换频、成功清除和失败清除次数依次为$n_{\mathrm{meas}}$、$n_{\mathrm{sw}}$、$n_{\mathrm{ok}}$和$n_{\mathrm{fail}}$，则总虚拟时间为
\begin{equation}
T_{\mathrm{total}}=\frac{D_{\mathrm{move}}}{5}+5n_{\mathrm{meas}}+n_{\mathrm{sw}}
                     +5n_{\mathrm{ok}}+3n_{\mathrm{fail}}.
\label{eq:q3-cost}
\end{equation}
成功清除的5秒包括光学定位3秒与清除2秒；清除不改变测向频道，无信号检测与末尾排查也须计费，任务结束后无需返回原点。

沿用全局符号，以$K_{j,t}$表示第$t$次更新后频道$j$的位置可行区域，$p_t$为机器狗当前位置，$\mathcal D_t$、$\mathcal C_t$分别为已发现和已清除频道集合，$\mathcal J_t$为尚未发现且未被完整排除的频道集合。真实位置$g_j$及源数$N$不进入决策。本问沿用第一问清除判据与第二问信息评价思路，以整场时间重新安排任务。


\subsubsection{从文献选点到多源任务安排}

为检验“单源选点有效，是否就能使整场任务更快”，本文回溯构造文献03的\textbf{谨慎贪心适配基线}。原算法沿估计协方差主轴的法向选取两个对称测点，随不确定性缩小而缩短与估计位置的距离，并选择移动较近的一侧\cite{VanderHook2014}。它提供了主动选点依据，但其高斯噪声、方位歧义与单目标设定不同于本题，不能直接套用滤波器及性能保证。

适配时，以角域交集的均匀面积矩确定中心、主轴与尺度，按原文选点规则取谨慎参数$\beta_{\mathrm c}=0.1$；将步长限制在可证明接收的范围内，并避免重复测点。均匀面积矩只用于选点，不作为真实后验分布；清除仍须通过19.5米最远距离判据。搜索使用原点及半径1200米圆周六点，优先就近完成已发现源，再选择覆盖收益较高且较近的扫描点，清除后不强制返回。两种策略采用相同的逐频道整单元覆盖停止标准。基线为本次补充验证实现，演练前冻结，并非原论文完整复现。

基线将搜索与各源定位分开安排：一个源的横向观测只服务自身，定位完成后的落点也未与后续扫描共同优化。我们此前的v4已引入扫描位置调整、共享站点和局部测向时间评价，但主要在“下一测点如何选”上节省费用。v6进一步比较完整任务的进入与离开位置，在长行程中补测其他目标，并把继续测向与提前试清放到同一时间尺度上评价。

\subsection{由观测约束到可靠清除}

\subsubsection{融合测向、接收与清除反馈}

对全向源，正常测向给出正向角域$W_i$与1500米接收上界；强信号给出$r_{\mathrm{strong}}$距离约束；无信号至少排除1000米范围。在检测点$s_i$取得反馈后，分别更新为
\begin{equation}
\begin{aligned}
 K_{j,t+1}&=K_{j,t}\cap W_i\cap B(s_i,1500), &&\text{正常示向度},\\
 K_{j,t+1}&=K_{j,t}\cap B(s_i,r_{\mathrm{strong}}), &&\text{强信号},\\
 K_{j,t+1}&=K_{j,t}\setminus B(s_i,1000), &&\text{无信号}.
\end{aligned}
\label{eq:q3-update}
\end{equation}
其中$B(s,r)$为闭圆盘。同一源的接收半径固定，若$s_+$有信号而$s_-$无信号，真实位置还满足
\begin{equation}
\|g_j-s_+\|\le R_j<\|g_j-s_-\|
\quad\Longrightarrow\quad
2(s_--s_+)^{\mathsf T}g_j<\|s_-\|^2-\|s_+\|^2,
\label{eq:q3-coupling}
\end{equation}
从而获得一个连接正、负观测的半平面约束。

实现使用相容位置的保守外包，以闭包计算最远距离。为覆盖$0.01^\circ$读数舍入，计算角域外扩至$1.0051^\circ$，题设误差界仍为$\alpha=\pi/180$。圆域采用外包多边形；普通负观测排除后可取凸包，光学失败产生的孔洞或分裂则保留，避免重复估计已排除区域的收益。

\subsubsection{根据区域尺度选择行动}

第一问表明，$d_{\max}(q;K_{j,t})\le r_{\mathrm{opt}}$即可保证清除。实现取$r_{\mathrm{safe}}=19.5$米：当$\rho(K_{j,t})\le19.5$米时，从$c(K_{j,t})$朝$p_t$的线段上，选择满足19.5米最远距离约束且离$p_t$最近的位置。该线段搜索减少最后一段移动，是第一问全域最近清除点的计算简化。

若区域半径超过120米，沿区域最长方向的法向，在圆心两侧各生成一个检测点，侧移距离取150米与$0.25\rho(K_{j,t})$中的较小值；区域较小时优先考虑中心。候选须满足接收的充分条件，再选移动较近者。例如，整个区域距候选点均不超过1000米即可保证接收；首次正观测也可用于构造接收下界。重复位置改用可接收侧移点。

\subsection{以逐频道覆盖保证搜索完整}

对未发现频道$j$，记实际无信号检测位置集合为$\mathcal P^-_{j,t}$，其尚未排除区域为
\begin{equation}
\mathcal U_{j,t}=\Omega\setminus
 \bigcup_{s\in\mathcal P^-_{j,t}}B(s,1000).
\label{eq:q3-uncovered}
\end{equation}
只有$\mathcal U_{j,t}$为空，才有依据认定该频道不存在源。经过某点但未检测该频道，或只检测其他频道，均不能计入覆盖记录。

计算时，将包围$\Omega$的正方形划分为边长28.125米的网格，保留所有与圆域相交的单元，包括边界单元。\textbf{只有单元四角均距该频道的同一个实际负观测点不超过999.99米，才标记整单元已排除。}圆盘为凸集，四角位于其中足以保证整个矩形被覆盖，略小于1000米的半径提供数值裕量；这也避免了仅检查网格中心造成漏判。

因此，正常停止须同时满足“已发现源全部清除”和“每个未发现频道的全部单元均已排除”；若已清除16个不同源，则可利用题设数量上限直接结束。计划覆盖和连续若干次无信号均不代替实际完成证据。该条件证明正常终止时没有漏源，不保证程序一定能够执行到终止。

\subsection{联合安排扫描与目标处理任务}

Engin和Isler对多个目标信息需求的联合考虑\cite{Engin2020}，启发本文让移动兼顾覆盖与定位。针对未知频道与逐频道计费，再以几何和路线估价实现任务安排；不移用文献成绩。

起点先扫描20个频道，再根据已发现目标方向设置250米侧向基线并筛选补测。在原点和半径1200米圆周六点的基础上，比较$0^\circ$至$50^\circ$、间隔$10^\circ$的六种旋转布局；目标的预计处理位置也可作为共享扫描位置。依据实际覆盖记录删除冗余扫描任务，并在保留必需网格角点覆盖约束的范围内调整站点，缩短相邻任务间路程。

扫描任务在站点执行；目标任务可能先测向，再移动到清除点。记任务$b$的预计离开位置为$e_b$：扫描任务取站点，目标任务以$c(K_{j,t})$近似。其进入位置$q_b(e_a)$由前一任务离开位置$e_a$决定，取可行清除点或下一测向点。候选序列$\pi$以当前位置为起点，按
\begin{equation}
\widehat T_{\mathrm{plan}}(\pi)=
\frac15\sum_{(a,b)\in\pi}
\left[\|q_b(e_a)-e_a\|+\|e_b-q_b(e_a)\|\right]
+6\,\#\mathcal J_t\,m_{\mathrm{scan}}(\pi)
\label{eq:q3-route-score}
\end{equation}
比较，其中$m_{\mathrm{scan}}$为扫描站数，6秒是一次测量与换频的保守估价。进入位置随前序任务变化，故路程代理通常不对称。通过多起点近邻构造和路径片段反转比较序列，只执行首项任务，再据反馈重规划；实际费用仍按式\eqref{eq:q3-cost}累计。

长移动还可提供补测机会。行程至少250米时，检查其四分之一、二分之一和四分之三处：目标半径至少100米，候选点距已有测点及最近尝试点均至少150米，距目标圆心不超过1000米；以指向圆心的方位及$\pm1^\circ$构造三种读数，预计半径平均下降至少100米才补测。多个位置可用时，比较各入选频道的预计半径下降减30米后的总和。站点补测还要求历史观测少于100次、距圆心不超过1200米且与历史测点及最近尝试点相隔至少150米。已有两次观测且距圆心超过100米时，三种读数须均达到19.5米精度，或平均半径降至原值的40\%以内。所有补测均逐频道计费。

\subsection{继续测向与提前试清的局部时间比较}

当$19.5<\rho(K_{j,t})\le100$米时，试清有可能比继续测向更快，但不能只比较失败3秒与测量5秒，还须考虑移动和失败后的后续定位。算法从当前位置、区域中心及邻近几何位置生成少量试清点，剔除重复或覆盖面积过小者。将区域代表位置与$-\alpha,0,\alpha$误差组合为有限情形集合$\Xi_{j,t}$。对情形$\xi=(x,e)$，试清点$q$的局部完成时间估计为
\begin{equation}
\widehat T_{\mathrm{try}}(q;\xi)=\frac{\|q-p_t\|}{5}+
\begin{cases}
5,&\|x-q\|\le r_{\mathrm{opt}},\\
3+\widehat T_{\mathrm{cont}}(K^-_{j,t}(q),q;\xi),&\|x-q\|>r_{\mathrm{opt}}.
\end{cases}
\label{eq:q3-optical-value}
\end{equation}
其中$K^-_{j,t}(q)$为失败后区域。后续过程最多试算五次测向，未完成则追加1000秒。动作按“80\%平均时间加20\%最大时间”评分；试清评分至少低1秒才采用，每源最多两次。这不是精确后验期望。

实际试清失败后，源距$q$大于20米。从可行区域中去除半径19.999米圆的内接多边形，保守保留非凸或分裂结构，再继续选择动作。图\ref{fig:q3-optical}中，试清前区域半径87.43米，试清与继续测向评分分别为24.90、29.59秒。失败后区域分裂，半径仍约87.43米；随后一次测向使其降至18.28米，再移至最远距离19.5米的位置成功清除。面积减少未必立即降低包围圆半径，失败反馈仍可用于后续定位。

\begin{figure}[htbp]\centering
\includegraphics[width=\linewidth]{q3-optical.pdf}
\caption{实际试清失败与后续定位。主图保持等比例，插窗为明确标注的局部放大。蓝色表示观测位置域，绿色表示试清失败后保留部分，灰色为后续测向前区域；圆圈分别为20米试清和19.5米可靠清除范围。局部坐标仅由平移、旋转得到，未使用真实源坐标。}
\label{fig:q3-optical}\end{figure}

\subsection{测试结果与策略分析}


\subsubsection{从基线缺陷到实际时间收益}

文献适配基线新增30场官方演练，v4与v6使用此前冻结的各50场记录。各场先以总虚拟时间除以实际源数，再对场次平均。全部结果均保留并核验官方日志，三组均全清，但地图与批次不同，不能视为同场配对。

\begin{table}[htbp]\centering\small
\caption{文献适配基线、早期v4与定稿v6的官方演练}\label{tab:q3-performance}
\begin{tabularx}{\linewidth}{L{2.7cm}Z Z Z Z Z}\toprule
策略 & 全清场数 & 源总数 & 均值 & 标准差 & 最大值\\\midrule
文献适配基线 & 30/30 & 373 & 333.59 & 37.21 & 396.52\\
v4 & 50/50 & 640 & 238.44 & 35.21 & 302.30\\
v6 & 50/50 & 645 & 231.71 & 31.20 & 303.41\\\bottomrule
\end{tabularx}
\par\vspace{3pt}\footnotesize 后三列单位为s/源；仅比较本题实现，不代表原文献算法排名。
\end{table}

相对文献适配基线，v6平均少用101.88秒/源（30.54\%），节省量的独立样本Welch 95\%区间为$[85.63,118.13]$秒/源。图\ref{fig:q3-performance}进一步显示，移动费用由277.35降至170.72秒/源，而检测费用由43.40增至46.70秒/源。\textbf{收益主要来自减少行走，而不是减少所有检测：适量补测使后续定位和扫描能顺路完成。}按两组共同源数层的合并频数标准化后，均值仍为329.64与234.74秒/源；这控制了源数构成，仍未消除地图、接收半径与批次差异。

\begin{figure}[htbp]\centering
\includegraphics[width=\linewidth]{q3-performance.pdf}
\caption{官方演练的逐场耗时与费用分解。左图保留30、50、50场观测，菱形与误差线为均值及其95\%置信区间；右图各分项先除以本场源数，再对场次平均。独立案例比较不单独识别某一模块的因果收益。}
\label{fig:q3-performance}\end{figure}

与已有协同机制的v4相比，v6仅再减少6.72秒/源（2.82\%），其中移动少6.22秒/源；节省量区间$[-6.48,19.93]$包含零，尚不能确认这一较小差距的稳定性。另有独立30场v6演练全部清除，均值226.34秒/源。因此，证据较明确地支持相对于逐源适配基线的整场改进，对v4则只观察到小幅改善，不能把二者写成同等强度的结论。

\subsubsection{三次正式测试记录}

按实际运行者的更正确认，\textbf{三次正式测试提交策略记为v6}，成绩见表\ref{tab:q3-formal}。自动归档的版本字段存在冲突，版本归属以该确认为依据；下表核验成绩与时间，不作为源码一致性或策略提速证明。正式元数据未给出真实源总数，故仅报告清除数，不另推算全清率。

\begin{table}[htbp]\centering\small
\caption{第三问已提交的正式测试记录}\label{tab:q3-formal}
\begin{tabularx}{\linewidth}{L{5.5cm}Z Z Z}\toprule
测试案例编码 & 清除源数 & 平均定位清除时间/s & 程序运行时间/s\\\midrule
BXUJ-W9QZ-KQ54-NKCU & 15 & 195.399 & 8.619\\
QNXC-8U8P-SQBF-A2A6 & 12 & 276.748 & 6.149\\
89UC-2392-TZ5M-ZUTB & 10 & 248.719 & 6.076\\\bottomrule
\end{tabularx}
\end{table}

平均定位清除时间包含搜索与失败操作；程序时间为官方接受进入至退出的真实时间差，不能与虚拟时间混用。原名日志及时间戳已核验。


\subsubsection{同一场景下，时间是怎样节省的}

为排除“只是地图更容易”的解释，另以统一校准本地模拟器，在30个预先固定的相同场景中分别运行两种冻结策略。源位置、接收半径和同位置固定误差场逐场一致，策略只读取动作反馈；这些是\textbf{本地配对实验，不是官方成绩}。两组均30/30全清，场均每源耗时由328.54降至231.89秒，30组均有下降；配对平均节省96.65秒/源，95\%区间为$[87.67,105.63]$。结论仍限于该近似分布。

图\ref{fig:q3-route}按“节省量最接近配对中位数”选取场景，避免挑选最大收益案例。相同16个源下，基线逐源横向定位并另赴固定扫描点；v6让扫描点随服务路线调整，将检测与清除穿插在连续行程中。路程从20.87降至12.68千米，节省1639.11秒移动时间；虽然多检测12次、多换频23次、试清失败5次，共增加98秒，整场仍少用1541.11秒。由此可以看出，改进的关键是用较小的观测支出换取更大的行程节省。

\begin{figure}[htbp]\centering
\includegraphics[width=\linewidth]{q3-route-coverage.pdf}
\caption{同一场景的本地路线对照（种子2026092410）。坐标、比例与源真值相同，连线来自实际动作，箭头示方向；圆点为至少执行四次检测的位置。真值只用于事后绘图，不输入策略。}
\label{fig:q3-route}\end{figure}

上述演练和配对结果支持已测场景的可用性，但另有本地固定误差场历史案例因试清的假想续跑区域数值退化而异常退出，当时真实位置域仍有效。冻结v6尚有这一估价异常传播风险，不能将失败前缀计作快速完成。先验、误差场及参数的系统敏感性仍待检验；正常终止的完整性与样本全清，均不证明软件必然可靠或时间最优。
